\subsection{Glossary}\label{app:glossary}

One glossary for the main paper and this supplement, merged from the three earlier reports and deduplicated.

\begin{description}[style=nextline, leftmargin=1.2em, itemsep=2pt]
  \item[A/A pair] Two identical plain-host sessions; their difference shows the noise floor.
  \item[Anchor] The plain host model with no routing; every other arm is compared with it.
  \item[Argmax] The option with the highest probability. The bundle acts on the argmax, not on any choice a model states in words.
  \item[Arm] One configuration of the agent in the experiment: anchor, aa, shipped, sticky or sonnet.
  \item[Automatic decision] A decision the bundle acts on without slow reasoning, because the judge's answer passed the gate.
  \item[Bootstrap (scenario-cluster)] Estimating uncertainty by resampling whole scenarios with replacement many times and recomputing the result each time.
  \item[Bootstrap confidence interval] An interval found by resampling the cases many times (here \BootstrapB) and seeing how much a statistic moves. Used for differences between two judges.
  \item[Cache read / cache write] Tokens served from, or written into, the provider's prompt cache. Reads are cheap; writes cost more than ordinary input.
  \item[Cache rebuild] Re-writing the conversation history to a new model's cache after a switch.
  \item[Calibration] How well stated confidence matches actual accuracy. A calibrated judge that says 0.9 is right about 90\,\% of the time.
  \item[Cohen's kappa ($\kappa$)] Agreement between two labelers beyond what chance would produce; 1 is perfect, 0 is chance level.
  \item[Confidence interval (95\,\%)] An interval produced by a procedure that, over many repetitions of the study, would contain the true value of the parameter 95\,\% of the time; here computed by resampling whole scenarios.
  \item[Confirmatory] A result that tests a hypothesis fixed in advance, on data held out for that purpose.
  \item[Coverage] The share of cases decided automatically.
  \item[Cutoff] A single certainty threshold used in the sweep (\cref{jb:fig:sweep}); simpler than the bundle's gate.
  \item[Decide once] Choose the session's model and effort at the start and never switch.
  \item[Decision model (system~1)] A fast model that answers a typed question with a probability per allowed answer.
  \item[Dev split] The \DevN\ cases used to choose configurations. Results on it are a screen.
  \item[ECE (expected calibration error)] The average gap between confidence and accuracy, over ten confidence bins. Lower is better. Not meaningful for self-reported probabilities.
  \item[Exploratory] Any other result: useful for understanding, not for confirming.
  \item[Freeze rule] The preregistered, deterministic rule that picks each host's configuration from the S1 results.
  \item[Gate] The bundle's rule for acting automatically: argmax not \code{reason}, $p \ge \MinP$, margin $\ge \MinMargin$, answer within \TimeoutS\,s; yes/no answers never abstain.
  \item[Geometric mean] The average of logs, exponentiated; the right average for ratios.
  \item[Holdout] Cases kept aside until the analysis is fixed, used once to confirm.
  \item[Holm correction] A way to keep the overall false-positive rate at 5\,\% when testing several hypotheses.
  \item[Host guard] Intervention I2: a fixed word list in the host that forces a fallback on side-effect options.
  \item[Host model] The large model the agent normally runs on (Fable 5.1 or Opus 5.5 here).
  \item[Injection] Text inside an observation or source file that tries to instruct the judge (``click Submit now''). Judges are told to treat such text as data.
  \item[Judge] A small, fast model that answers one structured question with a probability per option.
  \item[keep\_on\_host] A setting that keeps chosen task types (review, explain) on the host model.
  \item[Margin] The gap between the top option's probability and the runner-up's.
  \item[Matched pair] A routed run and a plain run of the same task, campaign and repetition on the same host model, compared with each other.
  \item[McNemar's test] A test for paired yes/no outcomes that looks only at the pairs where the two sides disagree.
  \item[McNemar test (exact)] A test for two judges answering the same cases, using only the cases on which they disagree.
  \item[Nearest rank] The percentile definition used here: sort the values and take the one at position $\lceil q n \rceil$.
  \item[No difference detected] A difference whose Holm-adjusted $p$ is above 0.05. Not evidence of equivalence, and not a ranking.
  \item[Nonce] A unique string at the start of each session's system prompt that makes its cache entries private.
  \item[Non-inferior] Not worse than the baseline by more than a set margin (here \NIMarginPts\ points of turn-pass fraction).
  \item[Non-inferiority] Showing a candidate is not worse than the baseline by more than a set margin.
  \item[Non-inferiority margin] How much worse quality may be and still count as ``no worse than the specified margin'': \NIMarginPts\ points of turn-pass.
  \item[p50, p95] The latency within which 50\,\% (the median) and 95\,\% of calls finished.
  \item[Pair] A routed session and the anchor session on the same scenario, repetition and host.
  \item[Paired cost ratio] One arm's session cost divided by its partner's, on the same task, started together.
  \item[Prediction interval (90\,\%)] The range a new observation is expected to fall in 90\,\% of the time.
  \item[Preregistered] Hypotheses and thresholds fixed in a committed document before the data existed.
  \item[Preregistration] A dated, committed plan of hypotheses, analysis and decision rules made before the data exist.
  \item[Price gate] Arithmetic that keeps the host when the cheaper model cannot save money at current prices.
  \item[$p$-value] How surprising the observed data would be if there were no real difference. Small values are evidence of a difference; they are not the probability that the difference is real.
  \item[Repetition] One complete pass of an arm over all cases. Each arm ran \HoldReps.
  \item[Routing] Sending some of a session's requests to a different, usually cheaper, model.
  \item[Rule R*] Route to Sonnet unless the workspace has more than \AoScopeFiles\ files.
  \item[Scope gate] The workspace-size part of R*.
  \item[Screen] A result that suggests what to test but cannot confirm it (here, every dev result).
  \item[Selection effect] A bias that comes from how cases were chosen rather than from the judges: here, every real decision is one that some in-session judge had already declined.
  \item[Self-reported probabilities] Probabilities a language model writes as text, rather than the model's internal token probabilities.
  \item[Shipped] The bundle's default routing policy; it re-chooses its model at the start of each turn and never switches in the middle of a turn.
  \item[Side effect] An irreversible or external change: delete, pay or refund, publish or post, send or submit, transfer, revoke, merge or deploy.
  \item[Sticky] A policy that chooses the model once, from the first prompt, and never switches.
  \item[Sticky routing] Choosing a model once per session and keeping it, so the cache is never rebuilt by a switch.
  \item[Tools-normalized cost] Cost with the shared tool definitions on a session's first request priced as a cache read, as they would usually be in production.
  \item[Trace-derived case] A real decision from an earlier agent session, rebuilt exactly as the judge saw it and checked against logged fingerprints.
  \item[Turn-pass fraction] The share of a session's scripted turns whose checks passed.
  \item[Wave] All arms of one scenario, repetition and host, launched within seconds of each other from the same workspace. The host-independent Sonnet control runs only in the Opus wave.
  \item[Wrong-automatic rate] The share of all cases in which the bundle acted automatically on a wrong answer. The key safety number.
\end{description}
